% Template per gli svolgimenti d'esame di Gasdinamica (prof. Favini).
% Uso: copiare in svolgimenti/<AAAA-MM-GG>_<ID o argomento>/svolgimento.tex,
% sostituire i segnaposto <<...>> e compilare con:  latexmk -pdf svolgimento.tex
% Le figure vanno in figure/ (cartella dello svolgimento), generate da calcoli.py.
% Ordine delle sezioni: CLAUDE.md, "Struttura dello svolgimento" (passi 1-11).
% Per gli esercizi con immagini, i passi 1-4 diventano: lettura del testo e
% dell'immagine, identificazione delle strutture, ipotesi ancorate all'immagine,
% misure sull'immagine; seguono ricostruzione del campo (passi 2-11), coerenza
% calcolo/immagine e figura annotata (vedi CLAUDE.md).

\documentclass[11pt,a4paper]{article}

% --- Preambolo di fonti/formulario.tex ---
\usepackage[T1]{fontenc}
\usepackage[utf8]{inputenc}
\usepackage[italian]{babel}
\usepackage{amsmath,amssymb}
\usepackage[a4paper,margin=2.5cm]{geometry}
\usepackage{array}

% --- Pacchetti aggiuntivi per gli svolgimenti ---
\usepackage{siunitx}
\usepackage{booktabs}
\usepackage{graphicx}
\usepackage{float}
\usepackage{hyperref}

\hypersetup{hidelinks}
\setlength{\parindent}{0pt}
\setlength{\parskip}{4pt}
\allowdisplaybreaks
\graphicspath{{figure/}}

% Numeri all'italiana: virgola decimale, unità SI.
\sisetup{output-decimal-marker={,}, group-separator={\,}, per-mode=symbol,
         exponent-product=\cdot, range-phrase={\text{--}}}

% --- Macro con la notazione del formulario (F) ---
% Il formulario usa direttamente \delta = (gamma-1)/2 (F §1) e q = (gamma-1)/(2 gamma) (F §8).
\newcommand{\Vv}{\boldsymbol V}                 % vettore velocità |V|
\newcommand{\sth}{s_{\rm th}}                    % entropia termodinamica
% Invarianti di Riemann 1D (F §3): R1 = a - delta u lungo u-a, R2 = a + delta u lungo u+a.
\newcommand{\Ru}{R_1}
\newcommand{\Rd}{R_2}
% Urti mobili (F §7): Mach relativo M_r = (u - w)/a, Mach dell'urto M_s = |M_{r1}|.
\newcommand{\Mr}{M_{r}}
\newcommand{\Ms}{M_{s}}
% Onde 1D e caratteristiche 2D (F §3, §10).
\newcommand{\uma}{u-a}                           % famiglia u-a
\newcommand{\upa}{u+a}                           % famiglia u+a
\newcommand{\Cp}{C^{+}}
\newcommand{\Cm}{C^{-}}
% Riferimento a una sezione del formulario o a una pagina del summary.
\newcommand{\F}[1]{(F~\S#1)}
\newcommand{\Sm}[1]{(S~p.~#1)}

\title{Gasdinamica -- Svolgimento \texttt{<<ID esercizio>>}}
\author{}
\date{Prova del <<GG/MM/AAAA>>}

\begin{document}
\maketitle

\section*{Traccia}
% Trascrizione fedele della traccia (o riassunto puntuale con tutti i dati e le richieste).
<<Testo della traccia.>>

\section{Lettura del testo e dati}
% Cosa chiede l'esercizio, dati con unità, grandezze richieste, eventuali
% interpretazioni della traccia (dichiarate e motivate).
<<Dati e richieste.>>

\begin{table}[H]
  \centering
  \begin{tabular}{lll}
    \toprule
    Grandezza & Valore & Note \\
    \midrule
    $\gamma$ & \num{1.4} & <<gas>> \\
    $R$      & \SI{287}{J/(kg.K)} & \\
    $p_1$    & \SI{1e5}{Pa} & <<stato 1>> \\
    \bottomrule
  \end{tabular}
  \caption{Dati del problema.}
\end{table}

\section{Modello e ipotesi}
% Modello scelto (1D instazionario, quasi-1D, 2D stazionario, ...) e ipotesi,
% ciascuna motivata: perché vale in questo problema.
<<Modello e ipotesi motivate.>>

\section{Analisi qualitativa}
% Previsione dei fenomeni prima del calcolo: quali onde, di che natura, in che
% verso; argomentazione fisica (non solo cosa, ma perché).
<<Analisi qualitativa.>>

\section{Topologia del campo}
% Struttura del campo d'onda: numerazione degli stati, onde che li separano,
% famiglia e natura di ciascuna onda (schema x--t o piano fisico).
<<Topologia, stati e onde.>>

\section{Bilancio equazioni/incognite}
% Se pertinente: incognite, equazioni disponibili, condizioni di raccordo.
<<Bilancio (oppure: non pertinente, perché ...).>>

\section{Scorciatoie fisiche}
% Se presenti: simmetrie, invarianti uniformi, stati noti a priori.
<<Scorciatoie (oppure: nessuna).>>

\section{Calcolo}
% Per ogni stato: formula analitica e risultato finale, SENZA i valori
% sostituiti; niente rimandi al formulario nel testo (solo nei commenti di
% calcoli.py). Per i procedimenti iterativi: primo tentativo come lo imposta il
% docente, criterio di convergenza, risultato finale. Niente codice.
\subsection{Stato <<n>>}
\[
  \frac{T_0}{T} = 1 + \delta M^2 = \num{1.8}.
\]

\section{Verifiche a posteriori}
% Ogni verifica scritta esplicitamente con il suo esito (verificata / non
% verificata e conseguenze). Se una fallisce: correzione dell'impostazione.
\begin{itemize}
  \item <<Verifica>>: <<esito>>.
\end{itemize}

\section{Risultati}
\begin{table}[H]
  \centering
  \begin{tabular}{lll}
    \toprule
    Grandezza & Valore & Esatto / approssimato \\
    \midrule
    <<grandezza>> & <<valore>> & <<esatto>> \\
    \bottomrule
  \end{tabular}
  \caption{Risultati richiesti.}
\end{table}

\section{Rappresentazioni}
% Diagrammi (x--t, p--u, p--theta, ...) generati da calcoli.py, assi e stati
% etichettati. Ogni figura va commentata nel testo.
% \begin{figure}[H]
%   \centering
%   \includegraphics[width=0.8\textwidth]{<<nome_figura>>}
%   \caption{<<Didascalia.>>}
%   \label{fig:<<nome>>}
% \end{figure}
<<Commento ai diagrammi.>>

\section{Discussione fisica}
% Significato dei risultati, limiti del modello, cosa cambierebbe variando i dati.
<<Discussione.>>

\end{document}
